% \section{Introduction}
% \label{sec:intro}

% \begin{figure}[t]
% \centering
% \includegraphics[width=0.95\linewidth]{landmark/fig1_hero_v6.pdf}
% \caption{\textbf{A converged classifier rejects the concept it already learned---in one view.}
% \emph{Top:} a language-model valence axis (\emph{V-axis}) orders the nine FACED
% emotions; the \emph{same} ordering appears in the cohort human EEG ($r{=}0.87$
% at emotion-category resolution; per-clip $r{=}0.71$) and, without ever being supervised on it, in a converged EEG
% classifier's class means ($|r|\in[0.60, 0.77]$ across ten checkpoints). The axis
% shown is the CLIP-text variant (the tightest brain match); the nine-story Qwen3
% probe of \S\ref{sec:vaxis-llms} gives $r{=}0.80$. Thin lines mark the axis positions.
% \emph{Bottom:} the converged classifier has already absorbed the concept, so
% adding the V-axis as an auxiliary loss \emph{lowers} FACED-9 balanced accuracy
% ($0.658{\to}0.629$, $\Delta{=}{-}0.029$, raw $p{=}0.004$; of $25$ variants, $0$ help
% and $10$ hurt at global Holm $p{<}0.05$), whereas averaging its ten checkpoints---which share the axis but
% differ in their seed-specific residual---sets a new state of the art ($0.6948$;
% best single model $0.6755$).}
% \label{fig:overview}
% \end{figure}

% Auxiliary losses that pull a network's features toward a known concept
% direction are a standard way to inject prior knowledge. We report a
% regime in which they systematically backfire. Take an EEG emotion
% classifier that has already converged, and add a loss aligning its
% features with a valence direction: across twenty-five such recipes on
% FACED~\citep{chen2023faced}, none improve balanced accuracy (BACC) and
% ten significantly reduce it (global Holm; sixteen uncorrected). The same loss that is within noise of zero on a
% weak backbone significantly hurts a converged one. We call this the \emph{saturation regularity}: once task
% labels alone have driven the network onto the target concept, further
% supervision can only deform an already-converged representation, without
% improving generalization. The finding cuts two ways: it explains the concept-alignment failures we report, and---by locating the recoverable gain in the seed-specific residual supervision cannot reach---it yields a new state of the art by ensembling rather than supervising (\Cref{fig:overview}).

% To make ``the concept the network has saturated on'' concrete, we extract it. Pass nine emotion-evocative stories (each
% expanded to paraphrases) through a language model, mean-pool the
% final-layer hidden states class by class, and take the first principal component; the resulting
% \emph{valence axis} (V-axis) predicts movie-review sentiment at AUC
% $0.832$ zero-shot, tracks the cohort-averaged EEG response of $123$
% subjects at the emotion-category level ($r{=}0.80$), and---in its
% CLIP-text form---is encoded, without supervision, by all ten
% same-architecture EEG checkpoints we tested.
% Recent work shows that neural networks encode human-interpretable
% concept directions in their internal
% features~\citep{kim2018tcav,arditi2024refusal}. We use the V-axis as a
% \emph{probe} of what brain-decoding networks encode, not as a training
% signal---which, as the saturation regularity shows, a converged network
% cannot benefit from.

% \textbf{A nine-story axis.} The nine stories are one per
% FACED~\citep{chen2023faced} class; the language model is Qwen3-4B,
% chosen as the lead because it sits at the centre of the alignment
% manifold in our $14$-LLM sweep, with the highest per-stimulus
% alignment to a behavioural valence reference (the recipe itself is
% model-agnostic, see Appendix~\ref{app:per-llm-deep-dive}). The SST-2
% benchmark grounds the AUC $0.832$ headline: with zero task labels the
% V-axis projection matches a supervised logistic regression on the same
% features trained on ${\approx}50$ labels ($0.821$ at $N{=}50$; a
% $5$k-label same-feature classifier reaches
% $0.924$)~\citep{socher2013recursive}. And the
% direction is not idiosyncratic to any one model: across the
% $14$ language models from $560$M to $32$B parameters, the V-axis is
% recovered across families, and within Qwen3 is essentially
% scale-invariant (off-diagonal $r{=}{+}0.995$ over $15$ size-pairs).

% \textbf{A valence ordering, consistent with the axis, in cohort EEG.}
% A ridge regression on $160$
% channel--band features (32 EEG channels $\times$ 5 frequency bands,
% differential-entropy) of the FACED cohort EEG predicts the V-axis position of
% each of the $28$ stimuli (video clips spanning the $9$ emotion classes); the regressor uses leave-one-stimulus-out
% cross-validation, no per-subject fitting, no auxiliary supervision.
% The Qwen3-4B V-axis reaches $r{=}0.80$ on this
% ridge and a CLIP-text emotion axis $r{=}0.87$, both at emotion-category
% resolution (nested-permutation $p<2{\times}10^{-4}$); a CLIP-text axis over
% the $28$ per-clip descriptions reaches $r{=}0.71$.
% The third claim --- that converged EEG classifiers encode the V-axis
% without ever being trained on it --- is quantified on the ten
% same-architecture SOTA-pool checkpoints trained on the $9$ emotion
% labels alone: each encodes the CLIP-text V-axis in its class-mean
% subspace ($|r|\in[0.60,0.77]$; the nine-story Qwen3 axis more weakly and
% per-checkpoint non-significantly, \S\ref{sec:convergence}).

% \textbf{Yet supervision fails.} The natural way to use such an
% alignment is an auxiliary loss pulling the EEG network's penultimate
% features toward the V-axis. Across $25$ such recipes---knowledge
% distillation~\citep{hinton2015distilling}, representational
% similarity~\citep{kriegeskorte2008rsa},
% contrastive~\citep{khosla2020supcon,vandenoord2018cpc},
% parameter-efficient fine-tuning (PEFT) adapters~\citep{hu2022lora}, and topographic
% losses---none help: ten produce significant decrements (global Holm; sixteen uncorrected), five families showing \emph{monotonic destruction}
% (Figure~\ref{fig:saturation-transition}). The effect depends on
% backbone strength (no measurable effect on the weak $d{=}3$ backbone, a
% significant negative on the converged $d{=}6$ SOTA backbone; a significant
% backbone$\times$loss interaction, $p{=}0.001$; on FACED depth and recipe
% covary with convergence, isolated by the CIFAR-10 epoch sweep). This is
% \emph{saturation}: distinct from overfitting, it is what happens once
% task labels alone---without concept or valence guidance---have driven the network onto the target direction,
% so added supervision moves the class-mean component further along the
% axis ($\Delta|r|$ up to ${+}0.36$) while test accuracy falls and
% training accuracy stays at ceiling.

% \textbf{Where the gain lives.} The same diagnosis locates the
% recoverable gain in the within-class residual that supervision cannot
% reach: the class-mean subspace (the ``basin'') is shared across seeds,
% while the residual is seed-specific and load-bearing. Ensembling across seeds averages
% out residual noise without disturbing the basin---per-checkpoint
% residual encoding predicts leave-one-out contribution at $r{=}{+}0.74$,
% and a $10$-checkpoint ensemble reaches FACED-9 SOTA at
% $\mathbf{0.6948}$ (val-selected single $\mathbf{0.6755}$), a $+10.5\%$
% improvement over EMOD~\citep{chen2025emod}. Replacing the
% language-derived KD prototypes with random orthonormal vectors leaves
% BACC statistically equivalent (TOST $p{=}0.001$, $n{=}20$; $\Delta{=}{-}0.001$): the gain is the \emph{shape} of the loss, not the
% \emph{content} of the prototypes, so the V-axis is a probe of what the
% network saturated on, not the supervised signal.

% \paragraph{Contributions.} \textbf{(1)} The \emph{saturation
% regularity}: concept-aligned auxiliary supervision cannot improve a
% \emph{converged} classifier and tends to hurt it---established across
% $25$ recipes on FACED-9, for both valence and arousal, and reproduced
% on a CIFAR-10 classifier, with a benefit present only while the backbone
% is still weak (\S\ref{sec:saturation}); it is consistent on a second EEG
% dataset, SEED-V, where auxiliary supervision likewise yields no gain
% (\citealp{liu2022seedv}; Appendix~\ref{app:seedv-full}).
% \textbf{(2)} A bias--variance seed ensemble that turns this mechanism
% into a positive prescription and a new FACED-9 state of the art
% (\S\ref{sec:convergence},~\S\ref{sec:sota}).
% \textbf{(3)} A nine-story \emph{V-axis probe}: a one-dimensional valence
% direction extracted from nine language-model emotion stories, validated
% as a genuine semantic axis (text sentiment zero-shot, cohort-EEG
% prediction, and encoding by converged EEG classifiers) and used to
% diagnose what the network has saturated on
% (\S\ref{sec:vaxis-llms}--\S\ref{sec:vaxis-brain}).

% \paragraph{Scope.} Four boundaries scope the
% result. \textbf{(i)}~The V-axis is not privileged: a generic valence
% axis probes comparably (\S\ref{sec:vaxis-llms}), and inside the
% ensemble the language-derived KD prototypes are statistically equivalent
% to random orthonormal vectors (TOST $\varepsilon{=}0.01$, $p{=}0.001$,
% $n{=}20$; $\Delta{=}{-}0.001$ BACC, paired $p{=}0.72$).
% \textbf{(ii)}~We do not claim single-trial or within-subject decoding:
% the brain correlation is on the EEG response averaged across $123$
% subjects, and the within-subject cohort-transfer correlation is near
% zero ($r{=}0.03$; single-subject decoding is a separate result, \S\ref{sec:vaxis-brain}).
% \textbf{(iii)}~Alignment is not universally harmful---on a weak backbone
% the same loss is within seed noise ($+0.007$ BACC, $p{=}0.41$) and on a
% still-converging CIFAR-10 classifier it helps; saturation is specific to
% the \emph{converged} regime. \textbf{(iv)}~We do not present saturation
% as a proven theorem for all architectures and datasets: we establish it
% empirically on FACED-9 and find it consistent on SEED-V. Limitations are
% collected in \S\ref{sec:discussion}.


\section{Introduction}
\label{sec:intro}

\begin{figure}[t]
\centering
\includegraphics[width=0.95\linewidth]{landmark/fig1_final.pdf}
\caption[Concept alignment and predictive benefit]{\textbf{Concept alignment and predictive benefit can diverge.} \textbf{(a)} Standardized emotion-class positions in CLIP text, cohort EEG ridge readouts, and class-mean features from one classifier checkpoint. \textbf{(b)} Adding V-axis supervision lowers FACED balanced accuracy under a strong recipe. \textbf{(c)} SatEnsemble combines improved training with checkpoint probability averaging.}
\label{fig:overview}
\end{figure}

Concept-aligned supervision seeks to improve learning by shaping representations around task-relevant knowledge~\citep{koh2020concept}. For EEG learning, it offers a way to connect neural signals to interpretable attributes of the task. Valence and arousal, for example, have been used to guide EEG emotion representations~\citep{chen2025emod}. Yet the value of such supervision depends on what the receiving model already represents. A concept can be meaningful and readily decodable while offering little additional value as a training target. This makes the choice of auxiliary supervision a practical design problem. An objective can change a representation as intended without improving the prediction it is meant to support. We therefore ask: \emph{when does stronger concept alignment cease to improve prediction, and what information should learning exploit next?}

For EEG emotion recognition, we study valence---the positive--negative dimension of emotion---on the nine-class FACED benchmark~\citep{chen2023faced}. We construct language-derived directions, or \emph{V-axis probes}, from representations of emotion descriptions. The story-based probe predicts sentiment on SST-2~\citep{socher2013recursive} without task-specific training (\Cref{tab:cross-modal}). Both story-based and CLIP-text~\citep{radford2021clip} variants correspond to cohort-averaged EEG responses (\S\ref{sec:vaxis-brain}). A 14-model sweep recovers valence structure across several model families (Appendix~\ref{app:per-llm-deep-dive}). Rephrasing yields similar projections, supporting robustness to wording changes (Appendix~\ref{app:vaxis-protocol}). Paraphrase averaging improves sentiment transfer over a single description in the lead model (\Cref{tab:nshot-curve}). Within trained EEG classifiers, valence is already encoded in class-mean features (\S\ref{sec:convergence}). \Cref{fig:overview}a compares standardized emotion-class positions in CLIP text, cohort EEG ridge readouts, and one classifier checkpoint. These results establish a task-relevant, already represented concept for testing whether further supervision improves prediction.

We then test whether explicitly reinforcing an already represented concept improves prediction. We evaluate 25 auxiliary-supervision recipes on a standard EMOD~\citep{chen2025emod} receiving baseline, spanning distillation, representation matching, contrastive objectives, and topographic constraints. None yields a statistically significant improvement. Ten significantly reduce accuracy after global Holm correction (\Cref{tab:full-vaxis}). Follow-up interventions on a stronger recipe increase concept alignment (\S\ref{sec:saturation}), yet reduce predictive accuracy (\Cref{fig:overview}b). These experiments establish the \emph{Saturation Regularity} in EEG learning: additional concept supervision can strengthen alignment without detectable predictive gains, and can even reduce accuracy.

We next examine how this pattern depends on training conditions and whether auxiliary gains require meaningful targets. Valence and arousal objectives with no detectable effect under a weaker FACED recipe reduce accuracy under the stronger recipe (\Cref{fig:saturation-transition}b). Complementary experiments on CIFAR-10~\citep{krizhevsky2009cifar} compare training stages within one architecture. Early auxiliary gains are no longer statistically detectable at convergence, and aligned targets never significantly outperform random targets at any tested stage (\Cref{fig:saturation-transition}a). An intervention on SEED-V~\citep{liu2022seedv} provides a second EEG comparison, with only a small change under the tested recipe (Appendix~\ref{app:seedv-full}). These experiments characterize the dependence on training conditions, while the random-target controls separate auxiliary-objective gains from evidence for the value of target semantics.

The saturation results leave open what predictive information remains beyond class-mean concept alignment. Across the analyzed checkpoints, class-mean valence encoding has similar strength, motivating a closer examination of within-class variation. We introduce \emph{ResDiag}, a residual diagnostic framework for analyzing predictive information beyond class-mean alignment. ResDiag separates shared class structure from within-class variation. To test the predictive role of concept-related residuals, we use directional ablation~\citep{arditi2024refusal} while preserving the class-mean subspace. Accuracy decreases in eight of ten checkpoints (\S\ref{sec:convergence}). To examine the link to model complementarity, we measure leave-one-out ensemble contribution as the performance change upon removing each checkpoint. Residual encoding correlates with this contribution ($r=0.74$), with uncertainty assessed by seed-clustered bootstrap (Appendix~\ref{app:convergence-deep-dive}). Together, these analyses identify useful information beyond class means and motivate combining complementary predictions.

Motivated by the supervision controls and residual analyses, we develop \emph{SatEnsemble}, a saturation-informed ensemble learning framework. Augmentation, prototype regularization, and increased depth strengthen individual models. Uniform averaging of checkpoint softmax probabilities combines their complementary outputs (\Cref{fig:overview}c; \S\ref{sec:sota}).

SatEnsemble achieves state-of-the-art EEG emotion recognition on FACED. The ten-checkpoint ensemble reaches 69.48\% balanced accuracy, outperforming both published EMOD~\citep{chen2025emod} and our five-seed EMOD ensemble baseline. The improved training recipe also yields 67.55\% with a validation-selected single model (\Cref{tab:sota}). Recipe ablations support gains from both improved training and prediction averaging (\S\ref{sec:sota}). Replacing language-derived prototypes with random orthonormal targets preserves performance within the tested equivalence margin (\Cref{tab:sota}). This control shows that prototype regularization can benefit prediction without relying on target semantics. These results establish a practical strategy for improving EEG prediction when further alignment of an already represented concept no longer helps.

To conclude, our contributions are:
\begin{itemize}
\item We identify the \emph{Saturation Regularity} in EEG learning: stronger concept alignment does not necessarily translate into better prediction. Systematic experiments characterize its dependence on training conditions and distinguish the effects of auxiliary supervision from those of target semantics.
\item We introduce \emph{ResDiag}, a residual diagnostic framework for analyzing predictive information beyond class-mean alignment. Ablations show that within-class residuals support prediction, while ensemble analyses link encoding strength to ensemble contribution.
\item We develop \emph{SatEnsemble}, a saturation-informed ensemble learning framework combining improved training with prediction averaging. It achieves state-of-the-art performance on FACED, with ablations supporting the value of both components.
\end{itemize}
